\documentclass[11pt]{article}

\usepackage[margin=1in]{geometry}
\usepackage{amsmath,amssymb,amsthm}
\usepackage{booktabs}
\usepackage{hyperref}
\usepackage{microtype}
\hypersetup{
  hidelinks,
  pdftitle={Dynamic Welfare-Maximizing Group Counting with Count-Valued Outcomes},
  pdfauthor={Vladimir; Hector; Francisco Marmolejo-Cossio}
}

\newtheorem{theorem}{Theorem}
\newtheorem{proposition}{Proposition}
\newtheorem{lemma}{Lemma}
\newtheorem{corollary}{Corollary}
\newtheorem{definition}{Definition}
\newtheorem{example}{Example}

\title{\textbf{Dynamic Welfare-Maximizing Group Counting with Count-Valued Outcomes:}\\
A Separation, Tractable Regimes, and Empirical Evidence}
\author{Vladimir \and Hector \and Francisco Marmolejo-Coss\'io}
\date{Preliminary draft --- July 9, 2026}

\begin{document}
\maketitle

\begin{abstract}
A welfare-maximizing testing planner uses a limited budget of pooled tests to
certify healthy individuals. Prior work studies static designs and dynamic
policies with binary outcomes. We study a richer dynamic model in which a test
reports the exact number of infected samples in its pool. The paper has three
parts. First, we give a closed-form separation: in a homogeneous high-prevalence
population, the best static design earns $Buq$, whereas a dynamic count-valued
strategy searches $kG$ people and earns at least
$u[1-(1-q)^{kG}]$ under the same strict hard-clearing objective. Because the
comparison changes both adaptivity and outcome resolution, we state it only as
a joint separation.
Second, we give efficient exact posterior algorithms for structured histories:
componentwise inference for disconnected test hypergraphs, an
elementary-symmetric-polynomial algorithm for fully observed laminar layers,
and a linear-time forward--backward algorithm for an acyclic chain with
one-bit separators. Third, exact experiments show mean relative improvements
of $0.63\%$, $3.97\%$, and $5.07\%$ on three small-instance regimes. A separate
resolution experiment finds that the
three-level channel $\{0,1,\geq2\}$ captures $84.5\%$ of the full counting gain
on a selected $N=6$ instance. Together, these results identify when richer
outcomes improve welfare, when exact inference remains tractable, and where the
present analysis stops.
\end{abstract}

\section{Introduction}

Imagine a planner facing 32 people in a high-prevalence population. Each person
is healthy with probability $q=0.05$, a healthy certification is worth $u$, and
the planner has seven tests. A fixed binary design does best by testing seven
people individually, for expected welfare $7uq=0.35u$. A count-valued policy can
instead test two disjoint pools of 16. If either pool contains a healthy person,
four count-guided splits isolate one such person, and a final singleton test
certifies that person under the strict hard-clearing rule. This policy earns at
least
\[
  u\bigl[1-(1-0.05)^{32}\bigr] \approx 0.806u.
\]
The mechanism is a counting argument. Individual tests inspect seven possible
healthy people; count feedback turns the remaining tests into a search over 32.

This example motivates dynamic welfare-maximizing group counting. A planner has
a heterogeneous population, a test budget, and a physical pool-size limit.
Welfare is earned only for people placed in a pool observed to contain no
infected samples. Static welfare-maximizing pooled testing fixes every pool in
advance \cite{finster2023}. Dynamic binary testing chooses the next pool after
observing whether each earlier pool was all-clear \cite{lopez2026}. We ask what
becomes possible when the observation is the full count of infected samples.

Count-valued observations are established in quantitative group testing, where
the usual objective is recovery with as few tests as possible
\cite{karimi2019}. Semi-quantitative group testing similarly replaces a binary
result with a small number of qPCR-derived bins \cite{nambiar2023}. Our novelty
is not the count channel itself. It is the welfare objective under strict hard
clearing, the separation from fixed designs, and the use of structured
count-posterior inference inside a dynamic decision problem.

The example changes two ingredients at once: the policy becomes adaptive and
the outcome becomes more informative. It therefore proves a joint separation,
not a value-of-counting result in isolation. The unresolved benchmark is the
optimal dynamic binary policy on the same homogeneous family; solving it would
separate the gains from adaptivity and resolution.

We make three contributions.

\begin{enumerate}
  \item We prove a closed-form joint separation under strict hard clearing,
  including the final certification test that the reward definition requires.
  \item We give exact posterior algorithms for verified tractable structures:
  disconnected components, fully observed laminar layers, and an acyclic
  one-bit-separator chain. These algorithms expose the structural reason exact
  inference can scale without enumerating all $2^N$ latent profiles.
  \item We consolidate two empirical views. A hierarchy experiment measures the
  dynamic count-versus-binary gap, while a resolution curve measures how much
  of that gap is captured by intermediate count channels.
\end{enumerate}

We keep the main text focused on this separation--inference--evidence
connection. Results specific to dynamic binary testing are inherited from the
predecessor study \cite{lopez2026}; posterior complexity, exact fiber
enumeration, and unrestricted sampling on count fibers are discussed only in
Appendix~\ref{app:complexity} or left for future work.

\section{Model and Comparison Classes}

There are $N$ individuals. Individual $i$ has utility $u_i\geq0$ and an
independent latent infection state $Z_i\sim\mathrm{Bernoulli}(p_i)$. Write
$q_i=1-p_i$ for the probability that $i$ is healthy. A test may contain at most
$G$ individuals, and the planner may perform at most $B$ tests.

For a pool $t\subseteq[N]$, the count-valued outcome is
\[
  r(t,Z)=\sum_{i\in t} Z_i.
\]
The corresponding binary outcome is $\mathbf{1}[r(t,Z)>0]$. The count channel
refines the binary channel: the binary result can be recovered from the count,
but not conversely.

A policy produces a history
$h_B=((t_1,r_1),\ldots,(t_B,r_B))$. Under \emph{strict hard clearing}, welfare is
credited only after a person belongs to an actually tested zero-count pool:
\[
  U(h_B,Z)=\sum_{i=1}^N u_i
  \mathbf{1}\!\left[\exists j\leq B:\ i\in t_j,\ r_j=0\right].
\]
Deduction alone is not rewarded. A policy that infers a person is healthy must
still place that person in a zero-count test to collect the utility.

We compare three classes.

\begin{definition}[Static binary design]
A static binary design fixes all $B$ pools before any outcome is observed.
Because welfare depends only on zero versus nonzero outcomes, additional
positive-count resolution cannot change the value of a fixed design.
\end{definition}

\begin{definition}[Dynamic binary policy]
A dynamic binary policy chooses $t_j$ as a function of the previous pools and
binary outcomes $\mathbf{1}[r_\ell>0]$, $\ell<j$.
\end{definition}

\begin{definition}[Dynamic count-valued policy]
A dynamic count-valued policy chooses $t_j$ as a function of the previous pools
and exact outcomes $r_\ell$, $\ell<j$.
\end{definition}

Let $\mathcal U^{\mathrm{stat}}$, $\mathcal U^{\mathrm{dyn,bin}}$, and
$\mathcal U^{\mathrm{dyn,count}}$ denote the optimal expected welfare in these
classes. Since a dynamic policy can ignore its history and a count policy can
coarsen its observations,
\[
  \mathcal U^{\mathrm{stat}}
  \leq \mathcal U^{\mathrm{dyn,bin}}
  \leq \mathcal U^{\mathrm{dyn,count}}.
\]
The hierarchy is definitional. The size of either gap is a substantive
question.

\section{A Separation for Dynamic Count-Valued Testing}
\label{sec:separation}

We now formalize the example from the introduction. Consider a homogeneous
population with common utility $u>0$ and common healthy probability
$0<q<1/2$. The population is otherwise exchangeable.

\begin{lemma}[Static optimum]
\label{lem:static}
For the homogeneous instance above, if $N\geq B$, the optimal static design
under strict hard clearing has value
\[
  \mathcal U^{\mathrm{stat}}=Buq.
\]
This remains true when the fixed pools may overlap.
\end{lemma}

\begin{proof}
A pool of size $g$ is all-clear with probability $q^g$ and then credits $gu$,
so its expected contribution is $ugq^g$. For $g\geq1$ and $q<1/2$,
\[
  gq^{g-1}\leq \frac{g}{2^{g-1}}\leq1,
\]
and hence $ugq^g\leq uq$. Summing over $B$ tests gives an upper bound $Buq$ for
disjoint pools.

For overlapping pools, write the welfare as a sum over individuals and use a
union bound on the events that certify each individual. Reversing the order of
summation again bounds the expected welfare contributed by each test by
$ugq^g\leq uq$. Thus overlap cannot exceed $Buq$. Testing $B$ distinct
individuals attains the bound.
\end{proof}

\begin{theorem}[Joint static-binary versus dynamic-count separation]
\label{thm:separation}
Let $G$ be a power of two, let $0<q<1/2$, and let
\[
  B=k+\log_2G+1
\]
for an integer $k\geq1$. If $N\geq\max\{B,kG\}$, then
\[
  \mathcal U^{\mathrm{dyn,count}}
  \geq u\left[1-(1-q)^{kG}\right].
\]
Consequently, whenever
\[
  Bq<1-(1-q)^{kG},
\]
the dynamic count-valued policy strictly outperforms every static design.
\end{theorem}

\begin{proof}
Use the first $k$ tests on $k$ disjoint pools of size $G$. If every person in
those pools is infected, the construction earns nothing. Otherwise at least one
tested pool has count strictly below its size and therefore contains a healthy
person.

Maintain a block known to contain at least one healthy person. Split it into two
equal halves and test one half. If its infected count is smaller than its size,
keep that half. Otherwise it is entirely infected, so keep the other half.
After $\log_2G$ tests, the maintained block is a singleton known to be healthy.
The last test places this singleton in an observed zero-count pool, which is
necessary under strict hard clearing.

The policy certifies at least one healthy person exactly on the event that the
initial $kG$ people contain one. That event has probability
$1-(1-q)^{kG}$.
\end{proof}

For $(q,G,k)=(0.05,16,2)$, the budget is $B=7$. Lemma~\ref{lem:static}
gives $0.35u$, while Theorem~\ref{thm:separation} gives $0.806u$. The lower
bound counts only one certified person, although zero-count pools encountered
along the way may certify more.

For fixed $B$ and power-of-two $G$, strict clearing gives
$kG=(B-\log_2G-1)G$. Equivalently, $G=2^{B-k-1}$, so integer $k\geq1$ gives
$kG=k2^{B-k-1}$. The maximum is $2^{B-2}$, attained at $k=1$ and also at
$k=2$ whenever that choice is feasible. The optimization sharpens the
construction but is not needed for the separation.

\section{Efficient Exact Inference in Structured Regimes}
\label{sec:structured}

A dynamic policy must update beliefs after observing counts. Unrestricted
histories couple many latent states, but simple overlap structure produces small
separators. The algorithms below compute exact posterior marginals without
enumerating every global profile. They are inference algorithms; integrating
each representation into a globally optimal large-scale decision policy is a
separate problem.

\subsection{Disconnected components}

Form a bipartite incidence graph between individuals and previously tested
pools. If this graph has disconnected components, the posterior factorizes
across them because the prior is independent and every count constraint lies
inside one component.

\begin{proposition}[Componentwise exact inference]
If every tested connected component contains at most $c$ individuals, exact
posterior marginals can be computed in
\[
  O\!\left(N+\sum_C 2^{|C|}(|C|+|h_C|)\right),
\]
where $h_C$ is the set of count constraints in component $C$. Thus the
calculation is linear in the number of components for fixed $c$, rather than
exponential in the full population size.
\end{proposition}

\begin{proof}
The likelihood indicator factors by component, as does the independent prior.
Enumerate feasible profiles separately within each component and normalize
locally. Products of the local normalizers recover the global normalizer, while
a marginal depends only on its own component.
\end{proof}

This decomposition is effective when adaptive tests form many small
disconnected regions. It provides exact inference within each region, but does
not imply global optimality of a greedy pool-selection rule.

\subsection{Fully observed laminar layers}

Suppose the tested pools form a nested hierarchy and every relevant level has
been observed. Subtracting child counts from parent counts fixes an exact count
inside each disjoint layer. For example, if $B\subset A$, with counts two and
five, then $A\setminus B$ contains exactly three infected individuals.

Consider one layer $L$ of size $m$ with observed infected count $d$. For
$0<p_i<1$, define odds $w_i=p_i/(1-p_i)$ and let $e_d(w_L)$ be the elementary
symmetric polynomial of degree $d$. When $d\geq1$, conditioning on
$\sum_{i\in L}Z_i=d$ gives
\[
  \Pr(Z_i=1\mid {\textstyle\sum_{j\in L}} Z_j=d)
  =\frac{w_i\,e_{d-1}(w_{L\setminus\{i\}})}{e_d(w_L)}.
\]
When $d=0$, every marginal in the layer is zero. Deterministic priors are removed
and their contribution is subtracted from $d$ before applying the formula.
Prefix and suffix recursions compute all required symmetric polynomials and all
$m$ marginals in $O(m(d+1))$ time. Each laminar layer is then solved independently.

\begin{proposition}[Fully observed laminar inference]
When nested observations determine disjoint layers $L_1,\ldots,L_s$ and their
counts $d_1,\ldots,d_s$, exact posterior marginals are computable in
\[
  O\!\left(\sum_{\ell=1}^s |L_\ell|(d_\ell+1)\right)
\]
time and $O(\max_\ell |L_\ell|(d_\ell+1))$ working memory.
\end{proposition}

The proposition applies only when every layer count is determined. Partially
observed laminar trees require message passing over unresolved parent counts,
which is outside the present algorithmic scope.

\subsection{An acyclic one-bit-separator chain}

Now consider the overlapping chain
\[
  \{0,1,2\},\ \{2,3,4\},\ \{4,5,6\},\ldots.
\]
Adjacent factors share one person, so a message needs to remember only one
binary state. A forward message $\alpha_j(s)$ stores the unnormalized mass of
all assignments through pool $j$ whose right boundary has state $s\in\{0,1\}$.
The transition sums over the one new interior state subject to the observed
count. A backward pass yields every boundary and interior marginal.

\begin{proposition}[Exact chain inference]
For $m$ size-three pools arranged as above, all posterior marginals can be
computed exactly in $O(m)$ time and $O(m)$ memory.
\end{proposition}

This is an acyclic factor graph with a one-bit separator. Its primal graph has
treewidth two, so calling it simply a treewidth-one model would be ambiguous.
More generally, junction-tree inference is polynomial in instance size at
fixed separator width and exponential in that width and the relevant factor
domain \cite{lauritzen1988}. Our prototype implementation covers this chain
special case; a general bounded-treewidth solver remains future work.

\section{Empirical Evidence}
\label{sec:empirical}

The experiments answer two different questions. The hierarchy experiment
compares optimal dynamic binary and count-valued policies. The resolution curve
then asks how much outcome detail is needed to capture the count-valued gain.

\subsection{Dynamic binary versus dynamic count}

For each configuration, infection probabilities are drawn independently from
$U(0,1)$ and utilities from $\{1,2,3\}$. The exact solvers compute the complete
value hierarchy. Table~\ref{tab:hierarchy} reports the resulting sample means.
The percentage is computed instance by instance and then averaged:
$100(\mathcal U^{\mathrm{dyn,count}}-\mathcal U^{\mathrm{dyn,bin}})/
\mathcal U^{\mathrm{dyn,bin}}$. It is not the ratio of the two aggregate means.
The generator is \texttt{augmented/hierarchy\_experiment.py} with PRNG seed 42;
raw rows are stored in \texttt{hierarchy\_small.csv} and
\texttt{hierarchy\_n7.csv} under \texttt{results/hierarchy}.

\begin{table}[t]
\centering
\caption{Exact dynamic binary and count-valued welfare. The $N=3$ and $N=5$
rows use 200 seeded instances; the heavier $N=7$ row uses 40. Standard errors
in the last column are in percentage points.}
\label{tab:hierarchy}
\small
\begin{tabular}{rrrrrr}
\toprule
$N$ & $B$ & $G$ & $\overline{\mathcal U}^{\mathrm{dyn,bin}}$
& $\overline{\mathcal U}^{\mathrm{dyn,count}}$ & relative gain $\pm$ SE \\
\midrule
3 & 2 & 3 & 2.686 & 2.704 & $0.63\%\pm0.09$ \\
5 & 3 & 5 & 4.600 & 4.771 & $3.97\%\pm0.19$ \\
7 & 3 & 7 & 5.454 & 5.717 & $5.07\%\pm0.37$ \\
\bottomrule
\end{tabular}
\end{table}

The hierarchy inequalities hold on every instance. The count-valued optimum is
strictly larger in 93 of 200, 196 of 200, and 40 of 40 instances, respectively;
the median relative gain is zero in the smallest regime. Because $N$, $B$, and
$G$ vary together, these rows motivate rather than replace a controlled sweep
over the three parameters.

\subsection{Resolution as a mechanism experiment}

A truncation channel reports $\min(r,c)$. The endpoint $c=1$ is binary, while
$c\geq G$ is the full count. The sweep contains eight designed instances. We
report \texttt{mixta\_n6\_g4} because it is the only recorded curve with a
nontrivial interior fraction, not as an estimate of how often such a curve
occurs. Its parameters are
$p=(0.15,0.25,0.35,0.45,0.55,0.65)$,
$u=(1,2,3,1,2,3)$, $B=3$, and $G=4$. On this instance,
\[
  \mathrm{frac}(c)=
  \frac{\mathcal U_c-\mathcal U_{\mathrm{bin}}}
       {\mathcal U_{\mathrm{count}}-\mathcal U_{\mathrm{bin}}}.
\]
The exact values appear in Table~\ref{tab:resolution}.

\begin{table}[t]
\centering
\caption{Resolution curve on the selected $(N,B,G)=(6,3,4)$ instance.}
\label{tab:resolution}
\small
\begin{tabular}{lrr}
\toprule
channel & optimal welfare & fraction of count gain \\
\midrule
$\{0\}/\{\geq1\}$ & 4.967923 & 0.000 \\
$\{0\}/\{1\}/\{\geq2\}$ & 5.297456 & 0.845 \\
$\{0\}/\{1\}/\{2\}/\{\geq3\}$ & 5.357733 & 1.000 \\
full count & 5.357733 & 1.000 \\
\bottomrule
\end{tabular}
\end{table}

Thus the three-level channel captures $84.54\%$ of the full binary-to-count
difference on this instance. This is a mechanism result on a selected small
case, not a population average. Other designed instances saturate at the same or
an earlier truncation level.

\subsection{Structured-inference checks}

The notebook compares the prototype structured routines with exhaustive
enumeration on small instances and then evaluates the same routines where
global enumeration is impossible. Table~\ref{tab:structured} records those
single-instance checks, not a validation suite. Timings are single-run
measurements on one machine and are only illustrative.

\begin{table}[t]
\centering
\caption{Executed correctness and scale checks for structured exact inference.}
\label{tab:structured}
\small
\begin{tabular}{llrr}
\toprule
structure & instance & maximum marginal error & time \\
\midrule
laminar layers & $N=12$ versus brute force & $2.2\times10^{-16}$ & $<1$ ms \\
laminar layers & $N=6000$ & global enumeration infeasible & 13 ms \\
one-bit chain & $N=13$ versus brute force & $1.1\times10^{-16}$ & not recorded \\
one-bit chain & 200 pools, $N=401$ & global enumeration infeasible & 1.1 ms \\
\bottomrule
\end{tabular}
\end{table}

\section{Discussion and Limitations}

\paragraph{Two mechanisms remain coupled.}
The separation is between static binary and dynamic count-valued testing. It
does not yet quantify the optimal dynamic-binary middle cell on the same
homogeneous family. The empirical hierarchy includes that cell on finite
heterogeneous instances, but it does not solve the analytic family in
Theorem~\ref{thm:separation}.

\paragraph{The algorithms address inference rather than global policy optimization.}
The component, laminar-layer, and chain routines compute exact posterior
marginals in their stated regimes. These marginals can support greedy scoring
or a small-state dynamic program, whereas a scalable exact policy optimizer
must also control the future action space.

\paragraph{Exact counts are an idealization.}
A quantitative laboratory measurement is noisy and often reports a range or a
continuous signal rather than an integer infection count. The resolution
experiment gives a first bridge to coarser channels, but a calibrated observation
model and external data are required before making a screening recommendation.

\paragraph{The experiments are small and synthetic.}
The exact comparison is valuable because it separates policies without
approximation error. Its population sizes are nevertheless limited by exact
dynamic optimization. The increasing percentages in Table~\ref{tab:hierarchy}
should be treated as evidence for a mechanism, not as an estimate for a
particular deployment.

\section{Conclusion}

Count-valued outcomes can change the economics of a small testing budget. The
separation family makes the gain visible with a simple comparison: static tests
inspect $B$ people one at a time, while a count-guided policy searches many more
and certifies a healthy person when one is present. Structured posterior
algorithms show that this information can be used exactly when overlap has small
separators. The experiments then connect the two ends: exact counts improve the
dynamic optimum, and a coarse three-level channel can capture most of that gain
on a selected small instance.

\appendix

\section{Monotonicity under Resolution Refinement}
\label{app:resolution}

A quantizer $Q$ partitions the possible counts. We restrict attention to
quantizers that isolate count zero, because otherwise an observation cannot
certify an all-clear pool under strict hard clearing.

\begin{lemma}[Refinement never lowers optimal welfare]
If $Q'$ refines $Q$, then $\mathcal U_Q\leq\mathcal U_{Q'}$.
\end{lemma}

\begin{proof}
Run an optimal $Q$-policy under $Q'$. At every stage, deterministically coarsen
the $Q'$ observation to its $Q$ bin and choose the pool prescribed by the
$Q$-policy. This reproduces the same coarsened history and the same tested pools
on every latent profile. Since both channels isolate zero, the realized
hard-clearing reward is identical. The optimal $Q'$-policy can only do better.
\end{proof}

The argument is a sequential deterministic-garbling construction, related to
the comparison of experiments \cite{blackwell1951}. With $B=1$, no observation
can affect a later action, so every zero-isolating resolution has the same
optimal value. The recorded resolution data include this flat-horizon control.

\section{Posterior Complexity, Fibers, and Sampling}
\label{app:complexity}

An exact-count history imposes linear constraints $Az=r$ on a binary latent
profile $z$. The posterior normalizer is a weighted sum over the feasible fiber
\[
  \mathcal F(A,r)=\{z\in\{0,1\}^N:Az=r\}.
\]
This representation explains both the informational value and the computational
cost of count feedback. We do not establish a general complexity theorem here:
both the complexity of the weighted normalizer and that of an individual
posterior marginal remain open in this model.

A small enumeration remains useful for intuition. For pools
$\{0,1,2\}$ and $\{2,3,4\}$, both observed with exact count one, only five
profiles satisfy the two equations. Replacing the exact outcomes with two
positive binary outcomes leaves 25 profiles. The all-infected profile belongs
to the binary set but not to the exact-count fiber. The comparison isolates how
exact counts sharpen the posterior, complementing the welfare separation in the
main text.

The implementation uses exact enumeration inside small connected components and
a Metropolis-style fallback for larger components. Local swap moves alone are
not sufficient: they preserve total infection count and can leave a fiber
disconnected. Alternating-path proposals repair verified small counterexamples,
but the prototype sampler does not prove irreducibility on every binary fiber.
A general sampler claim requires a connecting move set or a scoped
connectivity theorem.

\small
\begin{thebibliography}{9}

\bibitem{finster2023}
S. Finster, M. Gonz\'alez Amador, E. Lock, F. Marmolejo-Coss\'io,
E. Micha, and A. D. Procaccia.
\newblock Welfare-Maximizing Pooled Testing.
\newblock arXiv:2206.10660, 2023.

\bibitem{lopez2026}
N. Lopez, F. Marmolejo-Coss\'io, J. R. Tello Ayala, and D. C. Parkes.
\newblock Dynamic Welfare-Maximizing Pooled Testing.
\newblock arXiv:2601.22419, 2026.

\bibitem{blackwell1951}
D. Blackwell.
\newblock Comparison of Experiments.
\newblock In \emph{Proceedings of the Second Berkeley Symposium on Mathematical
Statistics and Probability}, pp. 93--102, 1951.

\bibitem{karimi2019}
E. Karimi, F. Kazemi, A. Heidarzadeh, K. R. Narayanan, and A. Sprintson.
\newblock Non-adaptive Quantitative Group Testing Using Irregular Sparse Graph
Codes.
\newblock arXiv:1910.06845, 2019.

\bibitem{nambiar2023}
A. Nambiar, C. Pan, V. Rana, M. Cheraghchi, J. Ribeiro, S. Maslov, and
O. Milenkovic.
\newblock Semi-Quantitative Group Testing for Efficient and Accurate qPCR
Screening of Pathogens with a Wide Range of Loads.
\newblock arXiv:2307.16352, 2023.

\bibitem{lauritzen1988}
S. L. Lauritzen and D. J. Spiegelhalter.
\newblock Local Computations with Probabilities on Graphical Structures and
Their Application to Expert Systems.
\newblock \emph{Journal of the Royal Statistical Society, Series B},
50(2):157--194, 1988.

\end{thebibliography}

\end{document}
